\textbf{Решение задачи 2-3}

\vspace{1em}
\textbf{1. Определение эффективного напряжения шума на входе}

Из условия задачи известно, что при воздействии только шума постоянная составляющая на выходе линейного амплитудного детектора (АД) $\bar{U}_{\text{ш}} = 0.1 \text{ В}$, а коэффициент передачи $K_{\text{д}} = 0.8$.

Связь между постоянной составляющей на выходе и эффективным (среднеквадратическим) напряжением шума на входе $\sigma$ для линейного АД описывается формулой:
\begin{equation*}
    \bar{U}_{\text{ш}} = K_{\text{д}} \sigma \sqrt{\frac{\pi}{2}}
\end{equation*}

Выразим и вычислим $\sigma$:
\begin{equation*}
    \sigma = \frac{\bar{U}_{\text{ш}}}{K_{\text{д}} \sqrt{\frac{\pi}{2}}} = \frac{0.1}{0.8 \cdot \sqrt{1.5708}} \approx \frac{0.1}{0.8 \cdot 1.2533} \approx \frac{0.1}{1.0026} \approx 0.0997 \text{ В}
\end{equation*}

\vspace{1em}
\textbf{2. Расчет постоянной составляющей при подаче гармонического сигнала}

На вход дополнительно подается гармонический сигнал с амплитудой $U_m = 0.3 \text{ В}$. Найдем отношение сигнал/шум на входе детектора:
\begin{equation*}
    a = \frac{U_m}{\sigma} = \frac{0.3}{0.0997} \approx 3.01
\end{equation*}

Так как $a \ge 3$, мы находимся в режиме «сильного сигнала» (полезный сигнал значительно превышает уровень шума). В этом случае математическое ожидание огибающей суммы гармонического сигнала и узкополосного шума (распределение Райса) можно найти по асимптотической формуле:
\begin{equation*}
    \bar{U}_{\text{с+ш}} \approx K_{\text{д}} U_m \left( 1 + \frac{1}{2a^2} \right)
\end{equation*}

Подставляем известные значения:
\begin{equation*}
    \bar{U}_{\text{с+ш}} \approx 0.8 \cdot 0.3 \cdot \left( 1 + \frac{1}{2 \cdot 3^2} \right) = 0.24 \cdot \left( 1 + \frac{1}{18} \right) = 0.24 \cdot \frac{19}{18} \approx 0.253 \text{ В}
\end{equation*}

\vspace{1em}
\textbf{3. Анализ эффективного напряжения шума на выходе АД}

Эффективное напряжение шума (СКО флуктуаций огибающей) на выходе АД зависит от уровня входного сигнала:
\begin{itemize}
    \item \textbf{При отсутствии сигнала ($U_m = 0$):} Закон распределения огибающей — релеевский. Эффективное напряжение шума, как было показано в задаче 2-2, равно:
    \begin{equation*}
        \sigma_{\text{вых, 0}} = K_{\text{д}} \sigma \sqrt{\frac{4-\pi}{2}} \approx 0.8 \cdot 0.0997 \cdot 0.655 \approx 0.052 \text{ В}
    \end{equation*}
    \item \textbf{При сильном сигнале ($U_m \gg \sigma$):} Закон распределения огибающей стремится к нормальному, а флуктуации определяются только синфазной составляющей входного шума. Эффективное напряжение шума на выходе возрастает и стремится к асимптотическому пределу:
    \begin{equation*}
        \sigma_{\text{вых, }\infty} \approx K_{\text{д}} \sigma = 0.8 \cdot 0.0997 \approx 0.08 \text{ В}
    \end{equation*}
\end{itemize}

\vspace{1em}
\textbf{4. Графики зависимостей}

Ниже приведен код для построения графика зависимости постоянной составляющей ($\bar{U}$) и эффективного значения шума ($\sigma_{\text{вых}}$) от амплитуды входного сигнала ($U_m$).

\begin{figure}[h]
    \centering
    \begin{tikzpicture}[scale=1.5]
        % Отрисовка осей
        \draw[->, thick] (0,0) -- (5,0) node[right] {$U_m$, В};
        \draw[->, thick] (0,0) -- (0,4) node[above] {$U$, В};
        
        % Сетка и разметка осей
        \draw[step=1cm, gray, very thin, dashed] (0,0) grid (4.5,3.5);
        \foreach \x/\xtext in {1/0.1, 2/0.2, 3/0.3, 4/0.4}
            \draw (\x,1pt) -- (\x,-1pt) node[below] {\xtext};
        \foreach \y/\ytext in {1/0.1, 2/0.2, 3/0.3}
            \draw (1pt,\y) -- (-1pt,\y) node[left] {\ytext};
            
        \node[below left] at (0,0) {0};
        
        % График постоянной составляющей \bar{U}_{с+ш}
        % Для построения используется аппроксимация \sqrt{\bar{U}_ш^2 + (K_d U_m)^2}, 
        % которая хорошо описывает переход от 0.1 В к асимптоте 0.8 * U_m
        \draw[domain=0:4.5, smooth, variable=\x, blue, very thick] 
            plot ({\x}, {10*sqrt(0.1^2 + (0.8*\x/10)^2)});
            
        % Асимптота постоянной составляющей (K_d * U_m = 0.8 * U_m)
        \draw[domain=0:4.5, dashed, blue] plot ({\x}, {0.8*\x});
        
        % График эффективного напряжения шума \sigma_{вых}
        % Плавный переход от 0.052 В к 0.08 В
        \draw[domain=0:4.5, smooth, variable=\x, red, very thick] 
            plot ({\x}, {0.8 - 0.28*exp(-(\x)^2)});
            
        % Асимптота эффективного напряжения шума (0.08 В)
        \draw[dashed, red] (0,0.8) -- (4.5,0.8);

        % Точка по условию задачи (U_m = 0.3 В)
        \filldraw[black] (3, 2.53) circle (1.5pt) node[above left] {$0.253$ В};
        \draw[dotted, thick] (3,0) -- (3,2.53);
        
        % Легенда
        \node[blue, right] at (1.5, 2.2) {$\bar{U}_{\text{с+ш}}$};
        \node[red, right] at (2, 0.95) {$\sigma_{\text{вых}}$};
    \end{tikzpicture}
    \caption{Зависимости постоянной составляющей и эффективного напряжения шума на выходе АД от амплитуды сигнала}
\end{figure}

\textbf{Ответ:} Постоянная составляющая выходного напряжения составит $\approx 0.253 \text{ В}$.